% optimization/realization_aware_optimization.tex

\chapter{Realization-Aware Optimization}

% ============================================================
\section*{Motivation}
% ============================================================

Classical alignment optimization operates directly on analytical target
geometry.

Real railway tracks, however, differ from their analytical target
states.

The physically realized alignment is influenced by:
\begin{itemize}
\item construction procedures,
\item machine behaviour,
\item ballast and track mechanics,
\item local correction limits,
\item discretization,
\item and measurement feedback.
\end{itemize}

Thus, optimization should minimize:
\[
J\bigl(\mathcal R(\kappa,\phi)\bigr),
\]
rather than:
\[
J(\kappa,\phi).
\]

Within AIM, optimization is therefore interpreted as optimization of
expected realized operational behaviour rather than of purely analytical
geometry.

% ============================================================
\section{Realization Operator}
% ============================================================

\begin{definition}[Realization operator]
The realization operator is:
\[
\mathcal R :
\mathrm{pose3}_{\mathrm{target}}
\mapsto
\mathrm{pose3}_{\mathrm{real}}.
\]
\end{definition}

The operator represents the transformation from analytical target state
to physically realized railway state.

The realization operator may include:
\begin{itemize}
\item sampling,
\item machine action,
\item smoothing,
\item local correction kernels,
\item and structural response of the track.
\end{itemize}

% ============================================================
\section{Optimization on Realized States}
% ============================================================

Objective functions should be evaluated on realized states rather than
on analytical target states.

\begin{definition}[Realization-aware optimization]
The general optimization problem is:
\[
\min_{\kappa,\phi}
\;
J\bigl(
\mathcal R(\kappa,\phi)
\bigr).
\]
\end{definition}

This formulation explicitly integrates realization effects into the
optimization process.

% ============================================================
\section{Pose3 Optimization}
% ============================================================

Optimization operates on the coupled pose3 state:
\[
(\gamma,\theta,\phi).
\]

Curvature and cant must therefore be optimized jointly rather than
independently.

This reflects the dynamic coupling between geometry and vehicle
behaviour.

% ============================================================
\section{Dynamic Objectives}
% ============================================================

\subsection{Effective Lateral Acceleration}

\begin{definition}[Effective lateral acceleration]
The effective lateral acceleration is:
\[
a_{\mathrm{eff}}
=
v^2\kappa-g\sin\phi.
\]
\end{definition}

This quantity represents the dynamically relevant acceleration acting
on the vehicle.

% ------------------------------------------------------------

\subsection{Jerk}

\begin{definition}[Jerk]
The jerk is:
\[
j
=
v^3\kappa'
-
gv\cos\phi\,\phi'.
\]
\end{definition}

Jerk constraints strongly influence transition design and passenger
comfort.

% ------------------------------------------------------------

\subsection{Typical Objective Components}

Typical optimization objectives include:
\begin{itemize}
\item curvature smoothness,
\item cant smoothness,
\item jerk minimization,
\item cant deficiency control,
\item vehicle comfort,
\item energy efficiency,
\item realization robustness.
\end{itemize}

% ============================================================
\section{Realization Constraints}
% ============================================================

Optimization must respect realization constraints such as:
\begin{itemize}
\item machine limits,
\item tamping limits,
\item ballast behaviour,
\item local correction ranges,
\item and measurement resolution.
\end{itemize}

Thus, the feasible design space is constrained not only by geometry and
standards, but also by realizability.

% ============================================================
\section{Inverse Realization}
% ============================================================

The optimization problem may be interpreted as an inverse realization
problem.

\begin{definition}[Inverse realization formulation]
The objective is:
\[
\mathcal R(\kappa_{\mathrm{target}},\phi_{\mathrm{target}})
\approx
(\kappa_{\mathrm{desired}},\phi_{\mathrm{desired}}).
\]
\end{definition}

The analytical target state therefore acts as a pre-image under the
realization operator.

% ============================================================
\section{Frequency Interpretation}
% ============================================================

The realization operator acts approximately as a low-pass filter on
curvature and cant.

High-frequency components are attenuated during realization.

Optimization must therefore avoid introducing geometrically meaningless
high-frequency corrections that cannot survive realization.

% ============================================================
\section{Sparse and Hybrid Optimization}
% ============================================================

Sparse and hybrid models are particularly suitable for
realization-aware optimization.

Sparse structures provide:
\begin{itemize}
\item numerical stability,
\item low parameter dimension,
\item and interpretable geometry.
\end{itemize}

Hybrid models additionally permit:
\begin{itemize}
\item local corrections,
\item realization compensation,
\item and adaptive refinement.
\end{itemize}

% ============================================================
\section{Runtime Optimization}
% ============================================================

Optimization may operate continuously during:
\begin{itemize}
\item construction,
\item maintenance,
\item measurement feedback,
\item and operational monitoring.
\end{itemize}

Thus, optimization becomes a runtime process rather than a purely
offline planning step.

% ============================================================
\section{Vehicle-Space Interpretation}
% ============================================================

The physically relevant object is not necessarily the track centerline,
but the motion of the railway vehicle within space.

This includes:
\begin{itemize}
\item vehicle center of gravity,
\item bogie motion,
\item vehicle envelope,
\item platform interaction,
\item and clearance behaviour.
\end{itemize}

Such interpretations connect realization-aware optimization to
Schwerpunkt-oriented alignment concepts
\cite{Hasslinger2002Schwerpunkt}.

% ============================================================
\section{Ellipsoid-Aware Optimization}
% ============================================================

For large-scale or high-precision systems, optimization may be performed
directly on the Earth reference surface.

In this case:
\begin{itemize}
\item curvature becomes geodesic curvature,
\item frames become surface-attached frames,
\item and world-to-track projection depends on ellipsoid geometry.
\end{itemize}

% ============================================================
\section{Key Insight}
% ============================================================

\begin{proposition}
The relevant optimization target is not the analytical alignment itself,
but the expected realized operational behaviour.
\end{proposition}

This fundamentally changes the interpretation of railway alignment
optimization.

% ============================================================
\section{Connection to AIM}
% ============================================================

\begin{summary}
Realization-aware optimization extends classical alignment optimization
toward physically realizable railway states.

Within AIM, optimization is formulated on:
\begin{itemize}
\item pose3 states,
\item realization operators,
\item runtime projection structures,
\item and dynamically relevant vehicle-space behaviour.
\end{itemize}

This provides a unified framework linking geometry, realization,
dynamics, optimization, and operational interpretation.
\end{summary}
